% ====================================================================
% paper_body.tex -- compact PRL-style scientific narrative.
% The full technical and diagnostic record lives in the analysis note.
% ====================================================================

Differential neutrino cross sections require correcting observed event samples
for detector response.  Iterative Bayesian unfolding on a binned migration
matrix is widely used, but the response tensor grows rapidly when additional
observables are reported.  OmniFold instead learns event weights in an unbinned
feature space through classifier-based density-ratio estimation
\cite{Andreassen:2019cjw,Andreassen:2021zzk}.  This makes simultaneous
higher-dimensional unfolding possible without specifying a separate binned
response tensor for every projection.

We apply this strategy to MINERvA medium-energy forward-horn-current
$\nu_{\mu}$ data.  We first reproduce the published double-differential
charged-current inclusive cross section in transverse and longitudinal muon
momentum, $d^{2}\sigma/(d\pT\,d\pPar)$ \cite{MINERvA:2021owq}.  We then add
available hadronic energy $\Eavail$, three-momentum transfer $q_{3}$, and
hadronic invariant mass $W$ to obtain simultaneous three-, four-, and
five-observable central values.  Existing MINERvA measurements reach three
reported dimensions, including quasielastic-like measurements of muon kinematics
and available recoil energy in two beam configurations
\cite{MINERvA:2022qe,MINERvA:2026qeleme}. Their signal differs from the inclusive
signal used here; our methodological advance is
the unbinned simultaneous construction and its projection to multiple
lower-dimensional observables.

The analysis uses the twelve MINERvA medium-energy playlists, corresponding to
$1.057\times10^{21}$ protons on target and an average neutrino energy near
\SI{6}{GeV}.  The signal definition, fiducial volume, flux treatment
\cite{MINERvA:2016iqn}, and reconstruction-level selection follow
Ref.~\cite{MINERvA:2021owq}.  After selection, the sample contains $4.12$
million data events and $32.8$ million simulated signal events, generated with
GENIE~2.12.6 and the MINERvA Tune~v1.

Each OmniFold iteration consists of two classification steps.  A
reconstruction-level classifier distinguishes data from simulated events and
defines a likelihood-ratio weight.  The event pairing transfers that weight to
truth, including unmatched events, and a truth-level classifier propagates a
smooth reweighting to the next iteration.  We use five iterations of
depth-three gradient-boosted decision trees with 100 trees.  The production
central value uses exact-split scikit-learn gradient boosting, whereas the
finalized two-dimensional systematic-universe and Poisson-bootstrap ensembles
use LightGBM at matched nominal capacity; the central value and its covariance
are therefore not estimator-matched.
Cross sections are formed by histogramming the reweighted truth events in a
declared reporting binning and dividing by integrated flux, target count, and
bin width.

The two-dimensional result provides a stringent end-to-end validation.  Its
integrated cross section is \SI{\sigTwoD}{cm^2/nucleon}, compared with
\SI{\sigTwoDpaper}{cm^2/nucleon} in the published analysis, a ratio of
$\ratioTot$.  Of the 205 reported bins, \SI{\binsTen}{\percent} agree within
\SI{10}{\percent}; the standardized bin differences formed with the published
uncertainties have an RMS of $\pullRMS$.  The standalone covariance combines
matched-central-value systematic universes with statistical and training-seed
blocks.  Its median relative uncertainty is \SI{\uqMedian}{\percent}, compared
with \SI{\uqPaper}{\percent} for the publication.  Figure~\ref{fig:validation}
shows the principal one-dimensional projections.  Because both results derive
from the same data, this validates the estimator rather than comparing
independent measurements; no combined-covariance statistic is quoted.

\begin{figure*}[t]
  \centering
  \includegraphics[width=0.78\textwidth]{model_comp_projections}\\[-0.4em]
  \begin{minipage}[c]{0.31\textwidth}
    \centering
    \includegraphics[width=\linewidth]{paper_validation_residual}
  \end{minipage}\hspace{1.2em}
  \begin{minipage}[c]{0.30\textwidth}
    \centering\footnotesize
    \textbf{Median relative uncertainty}\\[0.35em]
    \begin{tabular}{@{}lr@{}}
      This work & \SI{\uqMedian}{\percent} \\
      Published & \SI{\uqPaper}{\percent}
    \end{tabular}\\[0.8em]
    \textbf{Published-$\sigma$ standardized residuals}\\[0.35em]
    mean $=\pullMean$, RMS $=\pullRMS$
  \end{minipage}
  \caption{End-to-end validation against the published binned result
  \cite{MINERvA:2021owq}.  Top: $\pT$ and $\pPar$ projections of the published
  result, the five-iteration OmniFold result, and the MINERvA Tune~v1.  The two
  measurements use the same data, so their overlay validates the estimator
  rather than comparing independent results.
  Bottom left: $(\mathrm{OmniFold}-\mathrm{published})/\sigma_{\mathrm{published}}$
  on the reported two-dimensional grid; it is a descriptive shared-data
  residual, not an independent-result pull.  Bottom right: the independently
  reconstructed and published median relative uncertainties.}
  \label{fig:validation}
\end{figure*}

For the higher-dimensional results, every central-value construction is tested
against its lower-dimensional predecessor.  The three-, four-, and
five-observable marginal normalizations agree at the approximately
$1$--$2\%$ level, and injected distortions in each newly added variable are
recovered in closure tests.  These closure tests reweight the simulated signal
and do not exercise the background subtraction.  End-to-end pseudo-experiments
that include the purity-based background subtraction reproduce the input at
nominal truth except in the highest-$W$ cells, which are biased by up to
\SI{4}{\percent} in simulation, about one quoted standard deviation in one
$(\Eavail,W)$ cell; this bias is not included in the quoted uncertainties.
% Source: docs/orchestration/OUTCOME-20260925-s5c-purity-background-bias-at-high-W.md
% (D1: EW5-EW35 -3.5% to -4.1%, EW41 +1.66%); size vs adopted sigma:
% docs/orchestration/state/s5c/d1/bias_vs_adopted.json (max |bias|/sigma 1.07, EW35). KNOWN_ISSUES 75.

The $\Eavail$ dimension provides an independent view of the established
low-recoil discrepancy.  GENIE, the MINERvA Tune~v1, NuWro, and GiBUU all
underpredict the unfolded central value at low $\Eavail$.  Pion final-state
interaction variations change $d\sigma/d\Eavail$ by less than one percent,
whereas the empirical Valencia two-particle--two-hole contribution fills the
quasielastic--$\Delta$ dip. These tests constrain the chosen pion-FSI variations,
not all nuclear-transport models. Quasielastic-like measurements show sensitivity
to proton and pion FSI~\cite{MINERvA:2026qeleme}; their different signal definition
precludes identifying that effect with the inclusive discrepancy here.
Adding $q_{3}$ resolves the momentum-transfer
dependence of this feature.  On the maximal common $(\Eavail,q_{3})$ grid, the
central value agrees to within \SI{10}{\percent} with the dedicated published
low-recoil measurement~\cite{MINERvA:2022incl}.

The fifth observable, $W$, separates the remaining discrepancy from the
low-mass enhancement.  Figure~\ref{fig:joint} shows the joint
$(\Eavail,W)$ projection relative to the MINERvA Tune~v1, the tuned GENIE~2.12.6 simulation from which the unfolding starts.  The positive central-value
difference is concentrated where both $\Eavail$ and $W$ are large, rather than
being distributed uniformly across either one-dimensional projection.  Neither
one-dimensional projection resolves this structure.

\begin{figure}[t]
  \centering
  \includegraphics[width=\columnwidth]{paper_joint_localization}
  \caption{Localization in the joint $(\Eavail,W)$ central value, on the
  $7\times6$ reporting grid with $\Eavail$ edges
  $[0,0.1,0.2,0.4,0.8,1.5,3.0,100]\,\si{GeV}$ and $W$ edges
  $[0,1.1,1.4,1.8,2.2,3.0,100]\,\si{GeV}$; axes are cell indices counting from
  zero, so index~6 in $\Eavail$ and index~5 in $W$ are wide catch bins.  Left:
  the unfolded-to-Tune~v1 ratio per cell, dimensionless, on a $0.5$--$1.5$
  scale.  Right: the unfolded-minus-Tune~v1 difference as a
  \emph{cell-integrated} cross section in \si{cm^2/nucleon}, i.e.\ the
  density difference multiplied by that cell's
  $\Delta\Eavail\,\Delta W$; the catch-bin cells therefore carry their large
  widths and appear larger here than the ratio panel alone would suggest.  Read
  together, the panels put the remaining positive difference in the
  high-$\Eavail$, high-$W$ corner.  No significance is assigned.  The comparator is the MINERvA Tune~v1 (GENIE~2.12.6 with the analysis tune weights), not the untuned GENIE central value of Fig.~\ref{fig:generator-context}.}
  \label{fig:joint}
\end{figure}

The one-dimensional projections in Fig.~\ref{fig:generator-context} compare
the same central result with the four predictions that figure plots: GENIE's
own central value, the same GENIE with the Valencia two-particle--two-hole
component added, NuWro, and GiBUU.  All four underpredict the high-$W$ region.
Turning on the Valencia two-particle--two-hole component does not remove that
difference because it primarily populates lower $W$ --- visible here precisely
because the first two curves are two configurations of one generator rather
than independent predictions.  That several distinct generators agree on the
direction of the discrepancy is a useful model diagnostic.
Dedicated MINERvA shallow-inelastic measurements also report generator
disagreements~\cite{MINERvA:2025sis}, motivating differential tests across the
resonance--DIS transition. Their different selections preclude treating them as
independent confirmation of this localized difference.

\begin{figure*}[t]
  \centering
  \includegraphics[width=0.84\textwidth]{paper_eavailW_generators}
  \caption{Generator context for the unfolded $\Eavail$ (left) and $W$ (right)
  central-value projections.  GENIE-CV (untuned GENIE~2.12.10 on CH, with no two-particle--two-hole events), the same GENIE with the Valencia two-particle--
  two-hole contribution, NuWro~21.09, and GiBUU~2019 (both on carbon) all remain below the unfolded result
  at high $W$.  No significance is assigned.}
  \label{fig:generator-context}
\end{figure*}

These results demonstrate that an unbinned event-weight construction can be
validated against an established two-dimensional measurement and extended to
five simultaneous observables.  The higher-dimensional central values recover
the known low-recoil behavior and expose a localized high-$\Eavail$, high-$W$
generator deficit.  Every non-two-dimensional result in this Letter is a
central value.  A scalar five-dimensional covariance is adopted under a
digest-bound exception, and its \zprojCells-cell $(\Eavail,W)$ projection is
built; no superseded or historical covariance is used here.  We nevertheless
quote no publication-level significance, and the reason is measured rather
than precautionary.  In a dedicated campaign of two companion covariances built at two estimator-seed settings for this test (products distinct from the adopted bytes), varying the seed moves the projected
uncertainty by \SI{\sprojMeasured}{\percent}, against a
\SI{\sprojBound}{\percent} movement limit fixed before production.  That
movement did not fall with ensemble size over the range tested:
\SI{\seedEffectNforty}{\percent} averaged over four disjoint 40-throw subsets,
\SI{\seedEffectNeighty}{\percent} over the two 80-throw halves, and
\SI{\sprojMeasured}{\percent} on the full 160-throw ensemble, fitted exponent
$\seedEffectExponent$.  A same-seed resampling floor, measured at $N=40$ and
$N=80$, falls steeply over that step --- \SI{\floorNforty}{\percent} to
\SI{\floorNeighty}{\percent}, two-point exponent $\floorExponent$.
Statistical noise must fall with $N$; this did not fall between $N=40$ and
$N=160$, so the failing leg is not reporting the statistic's own resampling
noise.  The comparison is a single seed pair and the 40- and 80-throw points
are nested subsets of the same 160 throws, so the width of the seed-pair
distribution and the behaviour at much larger $N$ are both unmeasured; the
statistic is a maximum over the declared offset set, so further pairs can only
raise it.  Five of the covariance's 45 bands are
built at a fixed seed the variation cannot reach, carrying
\SI{\pinnedBandsSqrtTr}{\percent} of $\sqrt{\mathrm{Tr}\,C}$ and contributing
none of the movement, so their seed sensitivity is unprobed and the effect on
the total were they also to vary is not measured in either direction.  The same variation
moves the projected central values by at most
\SI{\cvMoveOverSigmaMax}{\percent} of their own quoted uncertainty (median
\SI{\cvMoveOverSigmaMed}{\percent}).

The selection-complete replacement underlying that covariance covers the five
detector bands implemented through shifted reconstructed kinematics.  MINOS
efficiency and the three GEANT hadron-interaction bands retain their
weight-based treatment.  This classification establishes their exclusion from
the kinematic replacement; it does not independently validate the completeness
of the hadronic-response model for the $\Eavail$ and $W$ measurement.

Two extensions lie outside this result.  Removing the published truth-level
muon cuts gives a preliminary extended-fiducial cross section in which
detector-supported and prior-dominated regions are reported separately.  A
full-event point-cloud classifier could replace the scalar learner while
preserving the two-step pairing; it remains method development because the
nominal/statistical pairing was declined, coverage has not been demonstrated,
and no point-cloud uncertainty product is adopted.  Both extensions are
documented in the companion analysis note.
